\section{武周的文字、宗教与政治传播} \label{sec:wu-zhou-literary-religious-politics}  \noindent	extbf{城市是交流的制度。} 武周的立后、封禅、建周、安西恢复、契丹危机与复唐，不能只按政治结果理解。都城、州城、港口、绿洲和军镇把翻译者、僧侣、使节、商人、工匠、书吏与军人放在同一条道路和市场之中；国家通过诏令、户籍、关津、寺院管理和外交礼仪试图整理这些流动。可是，城市遗址只能显示空间和物质活动，碑刻与文集多保存识字精英的声音，正史则倾向以朝廷的安全和秩序判断交流。它们需要互相校正，不能由任何一类材料单独推出整个城市或整个帝国的文化生活。  \noindent	extbf{宗教社群与文字传播。} 经卷、译场、寺院、道观、学校和科场使知识能够跨越政权和地域。它们既是信仰与学习的场所，也可能拥有土地、劳力、捐施和政治声望，因此会进入财政、法律和军政的计算。僧传的典范化、碑文的功德语言、诗文的修辞与官修史的褒贬都不是透明记录。较稳妥的写法，是确认某一文本、机构或政策在何时可见，同时保留社群人数、基层信仰与实际资产难以统计的事实。  \noindent	extbf{跨境材料。} “朝贡”“归附”“夷狄”等词首先是中原文书中的分类，不应直接等于稳定的统属或固定族群。突厥、吐蕃、回鹘、高句丽、新罗、日本以及中亚各方有自己的政治目标、书写传统和纪年方式；同一场战争、婚盟或使节往来在不同材料中可能有不同重点。将两类材料并读，不是要消除差异，而是防止朝廷的地理想象覆盖道路另一端行动者的选择。  \noindent	extbf{文化的物质条件。} 纸张、墨、颜料、金属、陶瓷、船只、马匹和粮食，决定了文本、图像和人能走多远。战争会破坏城市与寺院，也会造成新的迁徙、抄写、收藏和地方赞助；政权转换会改变年号和官职，却未必立即中断市场、亲属和宗教网络。考古和物质遗存可帮助追踪生产与交换的条件，但不能把器物的出土地点直接当作使用者身份或国家控制范围。文化史只有同交通、劳动和财政一起书写，才不会变成脱离社会的装饰。  \noindent	extbf{账本与限度。} 本节依托ST-655-WU-EMPRESS、ST-666-TAISHAN、ST-690-ZHOU-FOUNDATION、ST-692-ANXI-RECOVERY、ST-695-KHITAN-REBELLION、ST-705-TANG-RESTORATION等账本节点，并以两《唐书》、诏令会要、佛教文献、碑刻与边疆材料为主要来源入口。事件可锁定时间、行动与材料链，不能替代对文本体裁、保存地点、传本层次和区域覆盖的判断。我们可以说某城在某年易手、某僧侣或使节进入某一网络、某项礼仪或禁令被公布；不能据此判定所有居民的态度、所有社群的规模或某种交流的单一方向。证据的限度正是理解多元世界如何被叙述的前提。  